\documentclass[12pt]{article}
\usepackage[utf8]{inputenc}
\usepackage[T1]{fontenc}
\usepackage[margin=1in]{geometry}
\usepackage{setspace}
\usepackage{xcolor}
\usepackage[colorlinks=true, urlcolor=blue!60!black, linkcolor=blue!60!black]{hyperref}
\setlength{\parindent}{2em}
\setlength{\parskip}{1.25em}
\renewcommand{\baselinestretch}{1.0}
\begin{document}
\begin{center}
{\LARGE \textit{Martin v. Hunter's Lessee}, 14 U.S. (1 Wheat.) 304 (1816)\par}
\vspace{0.8em}
{\large Case Brief\par}
\end{center}
\vspace{1em}

\section{Memory Jogger}

In \textit{\href{https://www.courtlistener.com/opinion/85160/martin-v-hunters-lessee/}{Martin v. Hunter's Lessee}}, the Supreme Court held that Article~III and Section~25 of the Judiciary Act of 1789 constitutionally authorize the Supreme Court to exercise appellate review over final state-court civil judgments involving federal constitutional, statutory, or treaty questions to preserve national supremacy and uniformity of federal law.

\section{Detailed Case Facts}

Thomas, Sixth Lord Fairfax, a citizen and inhabitant of Virginia, held proprietary fee-simple title under seventeenth-century English royal letters patent (confirmed by Virginia statutes of 1736 and 1748) to the vast territory known as the Northern Neck of Virginia (\href{https://supreme.justia.com/cases/federal/us/14/304/}{14 U.S. at 307--11}). Lord Fairfax operated a proprietary land office in the Northern Neck from 1748 until his death in December 1781. By his last will and testament, he devised the Northern Neck lands in fee simple to his nephew, the Reverend Denny Martin---a native-born British subject who resided in England throughout the Revolutionary War and never became an American citizen---on condition that he assume the name and arms of Fairfax, which he did, becoming Denny Fairfax (\href{https://supreme.justia.com/cases/federal/us/14/304/}{14 U.S. at 311--12}).

During and after the American Revolution, Virginia enacted statutes restricting alien landownership and providing for the escheat and confiscation of Loyalist estates, though it was stipulated in the agreed record that the particular parcel of ``waste and ungranted'' land at issue had never been the subject of an inquest of office or formally escheated and seized into the hands of the Commonwealth (\href{https://supreme.justia.com/cases/federal/us/14/304/}{14 U.S. at 312}). On April 30, 1789, pursuant to a land-office treasury warrant issued on January 23, 1788, the Governor of Virginia issued a land patent granting the disputed parcel to David Hunter, a Virginia citizen, who entered into possession (\href{https://supreme.justia.com/cases/federal/us/14/304/}{14 U.S. at 312}).

Meanwhile, two federal treaties between the United States and Great Britain intervened to protect British subjects' landholdings: the Definitive Treaty of Peace of 1783 (which prohibited future confiscations of British property) and the Treaty of Amity, Commerce, and Navigation of 1794 (the Jay Treaty), which confirmed titles already held by British subjects (\href{https://supreme.justia.com/cases/federal/us/14/304/}{14 U.S. at 313, 360}). During the prolonged litigation, Denny Fairfax died (between 1796 and 1803), and his heir and devisee, Philip Martin, succeeded to his claim (\href{https://www.fjc.gov/history/cases/cases-that-shaped-the-federal-courts/martin-v-hunter-s-lessee}{Federal Judicial Center}; \href{https://supreme.justia.com/cases/federal/us/14/304/}{14 U.S. at 311--12}). Notably, Chief Justice John Marshall and a syndicate of relatives had contracted to purchase substantial portions of the Fairfax estate from Martin, creating a direct financial conflict of interest that required Marshall's recusal (\href{https://supreme.justia.com/cases/federal/us/14/304/}{Justia Overview}; \href{https://www.fjc.gov/history/cases/cases-that-shaped-the-federal-courts/martin-v-hunter-s-lessee}{FJC}).

\section{Procedural History}

\begin{enumerate}
    \item \textbf{Winchester District Court (Virginia, 1791--1794):} In April 1791, David Hunter's lessee brought an ejectment action in the Virginia District Court at Winchester against the tenants in possession; Denny Fairfax was admitted to defend title. On a stipulated statement of facts treated as a special verdict, the District Court entered judgment for the defendant Fairfax on April 24, 1794 (\href{https://supreme.justia.com/cases/federal/us/14/304/}{14 U.S. at 306}).
    \item \textbf{Virginia Court of Appeals (1810):} At its April 1810 Term, the Virginia Court of Appeals (then Virginia's highest court) reversed the District Court and entered judgment for Hunter (\href{https://supreme.justia.com/cases/federal/us/14/304/}{14 U.S. at 306--07}).
    \item \textbf{First U.S. Supreme Court Appeal---\textit{Fairfax's Devisee v. Hunter's Lessee}, 11 U.S. (7 Cranch) 603 (1813):} On writ of error under Section~25 of the Judiciary Act of 1789, the U.S. Supreme Court (in an opinion by Justice Story, with Marshall recused) reversed the Virginia Court of Appeals, holding that because Virginia had not completed an inquest of office before the Jay Treaty of 1794 took effect, the treaty protected Fairfax's title. The Court issued a mandate directing the Virginia Court of Appeals to execute its judgment (\href{https://www.fjc.gov/history/cases/cases-that-shaped-the-federal-courts/martin-v-hunter-s-lessee}{FJC}; \href{https://supreme.justia.com/cases/federal/us/14/304/}{14 U.S. at 305, 313}).
    \item \textbf{Virginia Court of Appeals Refusal---\textit{Hunter v. Martin}, 4 Munf. 1 (Va. 1815):} Led by Judge Spencer Roane and Judge William Cabell, the Virginia Court of Appeals unanimously refused to obey the Supreme Court's mandate, declaring that Article~III did not extend federal appellate jurisdiction to sovereign state courts, that Section~25 of the Judiciary Act of 1789 was unconstitutional insofar as it purported to do so, and that the Supreme Court's prior proceedings were \textit{coram non judice} (\href{https://supreme.justia.com/cases/federal/us/14/304/}{14 U.S. at 305--06}; \href{https://chicagounbound.uchicago.edu/cgi/viewcontent.cgi?article=1427&context=journal_articles}{Chicago Unbound}).
    \item \textbf{Second U.S. Supreme Court Appeal (1816):} Philip Martin brought a second writ of error under Section~25 directly challenging the Virginia Court of Appeals' refusal to execute the mandate. Argued March 12--14, 1816, and decided March 20, 1816, the Supreme Court reversed the Virginia Court of Appeals and directly affirmed the 1794 judgment of the Winchester District Court (\href{https://supreme.justia.com/cases/federal/us/14/304/}{14 U.S. at 304, 362}).
\end{enumerate}

\section{Judicial Votes}

The Court decided the case by a unanimous \textbf{6--0} vote among the participating Justices, with Chief Justice John Marshall recused (\href{https://supreme.justia.com/cases/federal/us/14/304/}{Justia}; \href{https://www.fjc.gov/history/cases/cases-that-shaped-the-federal-courts/martin-v-hunter-s-lessee}{FJC}):

\begin{itemize}
    \item \textbf{Majority Opinion (5 Justices joining in full):} Justice Joseph Story (author), joined by Justices Bushrod Washington, Henry Brockholst Livingston, Thomas Todd, and Gabriel Duvall.
    \item \textbf{Concurring Opinion (1 Justice):} Justice William Johnson concurred in the judgment and in the constitutionality of Section~25's revising power over cases and parties, while writing separately to reject several steps in Story's reasoning.
    \item \textbf{Dissenting:} None.
    \item \textbf{Recused / Not Participating (1 Justice):} Chief Justice John Marshall recused himself because of his personal and family financial interest in the Fairfax estate (\href{https://supreme.justia.com/cases/federal/us/14/304/}{Justia}).
\end{itemize}

\section{Holding}

\textbf{Judgment:} Reversed the judgment of the Virginia Court of Appeals rendered on the mandate, and affirmed the judgment of the Virginia District Court held at Winchester in favor of the Fairfax/Martin title (\href{https://supreme.justia.com/cases/federal/us/14/304/}{14 U.S. at 362}).

\textbf{Rule Announced:} Under Article~III, Section~2 and the Supremacy Clause of Article~VI, the appellate jurisdiction of the U.S. Supreme Court extends to final judgments of state courts in civil cases involving federal constitutional, statutory, or treaty questions, because Article~III attaches jurisdiction to the \textit{case} rather than the \textit{court} in which it arises. Consequently, Section~25 of the Judiciary Act of 1789 is constitutional, and when reviewing a state judgment denying a right or title claimed under a federal treaty, the Supreme Court may examine the underlying state-law title question on the face of the record to determine whether the treaty protects it (\href{https://supreme.justia.com/cases/federal/us/14/304/}{14 U.S. at 338, 351, 358--59}).

\textbf{Holding vs.\ Dicta:}
\begin{itemize}
    \item \textbf{Binding Holding:} (1) Section~25 of the Judiciary Act of 1789 is a constitutional exercise of Article~III appellate jurisdiction over final state-court judgments denying a federal right, title, or statutory validity; (2) a state court's refusal to obey a Supreme Court mandate is itself a final judgment reviewable on a second writ of error under Section~25; and (3) under Section~25, the Court may review a state-law title question intertwined with a federal treaty claim so that state courts cannot evade federal review by mischaracterizing the underlying title (\href{https://supreme.justia.com/cases/federal/us/14/304/}{14 U.S. at 351--59}).
    \item \textbf{Dicta and Reserved Questions:} Justice Story's elaborate argument that Article~III's ``shall be vested'' and ``shall extend'' language imposes a mandatory duty on Congress to vest the \textit{entire} federal judicial power---and specifically to create inferior federal courts to handle ``first-tier'' cases (federal question, ambassadorial, and admiralty cases) where Article~III uses the word ``all''---was unnecessary to resolve the validity of Section~25 and is classic dicta (\href{https://supreme.justia.com/cases/federal/us/14/304/}{14 U.S. at 328--37}; \href{https://constitution.congress.gov/browse/essay/artIII-S1-1/ALDE_00013512/}{Constitution Annotated}). In addition, the Court expressly reserved whether it possessed authority to issue a writ of mandamus directly to the Virginia Court of Appeals, avoiding a direct enforcement clash by affirming the Winchester District Court directly (\href{https://supreme.justia.com/cases/federal/us/14/304/}{14 U.S. at 362}).
\end{itemize}

\section{Analysis of Opinions}

\subsection{Majority Opinion (Justice Story)}

Justice Story anchored the Court's opinion in popular sovereignty rather than state-compact theory. Echoing the Preamble, he declared that the Constitution was ordained not by sovereign states, but by ``the people of the United States,'' who were fully competent to make state governmental powers subordinate to national authority in enumerated spheres (\href{https://supreme.justia.com/cases/federal/us/14/304/}{14 U.S. at 324--25}). Because a constitution must endure through ``a long lapse of ages,'' its grants of power are expressed in general terms and must receive a reasonable construction rather than one artificially narrowed to preserve pre-existing state prerogatives (\href{https://supreme.justia.com/cases/federal/us/14/304/}{14 U.S. at 326--27}).

Turning to Article~III, Story first advanced his famous mandatory-vesting and two-tier thesis: ``shall be vested'' and ``shall extend'' are imperative, and Article~III's repeated use of ``\textit{all} cases'' for federal-question, ambassadorial, and admiralty matters---contrasted with the omission of ``all'' for party-based ``controversies''---signaled that federal judicial power must extend in either original or appellate form to every case in the first tier (\href{https://supreme.justia.com/cases/federal/us/14/304/}{14 U.S. at 328--36}). Yet Story pivoted immediately to an independent textual ground: even if Congress had discretion whether to vest jurisdiction, Article~III's appellate grant contains no restriction as to the tribunal from which an appeal comes. ``It is the case, then, and not the court, that gives the jurisdiction'' (\href{https://supreme.justia.com/cases/federal/us/14/304/}{14 U.S. at 338}).

Story reinforced this textual reading with four structural and historical arguments:
\begin{enumerate}
    \item \textbf{Supremacy Clause Integration:} Article~VI explicitly contemplates that state judges will adjudicate federal questions arising incidentally in ordinary state-court litigation (such as a Contract Clause or Ex Post Facto defense). Because original jurisdiction has already lawfully attached in state court, federal judicial power can extend to ``all'' such cases only through appellate review (\href{https://supreme.justia.com/cases/federal/us/14/304/}{14 U.S. at 340--42}).
    \item \textbf{No Inconsistency with State Sovereignty:} The Constitution already restricts state sovereignty extensively (e.g., Article~I, Section~10), and federal courts routinely invalidate acts of state legislatures and executives; reviewing state judicial acts is no greater intrusion (\href{https://supreme.justia.com/cases/federal/us/14/304/}{14 U.S. at 343--44}). Moreover, under the Articles of Confederation, federal prize courts had exercised appellate review exclusively over state tribunals (\href{https://supreme.justia.com/cases/federal/us/14/304/}{14 U.S. at 345}).
    \item \textbf{Local Bias and National Uniformity:} The Constitution presumes that state attachments, prejudices, or interests might occasionally influence state courts, and---even assuming equal learning and integrity among state judges---without a single revising tribunal to harmonize ``jarring and discordant judgments,'' the Constitution, laws, and treaties of the United States would carry different meanings in different states (\href{https://supreme.justia.com/cases/federal/us/14/304/}{14 U.S. at 346--48}). Pre-judgment removal alone cannot solve this problem, because removal is itself a form of appellate jurisdiction and would be unworkable in many incidental-defense or criminal contexts (\href{https://supreme.justia.com/cases/federal/us/14/304/}{14 U.S. at 349--51}).
    \item \textbf{Contemporaneous Exposition and Practice:} Ratification debates, the First Congress's enactment of the Judiciary Act of 1789, and decades of unchallenged state-court acquiescence confirmed the constitutionality of Section~25 (\href{https://supreme.justia.com/cases/federal/us/14/304/}{14 U.S. at 351--52}).
\end{enumerate}

Finally, Story held that Virginia's refusal to obey the 1813 mandate was a final judgment denying the validity of a federal statute under Section~25, and---waiving finality objections to address the merits of the 1813 jurisdiction---reaffirmed under \textit{Smith v. Maryland}, 10 U.S. (6 Cranch) 286 (1810), that the Court must be able to examine the underlying state-law title on the record whenever a state court's denial of a treaty right rests on finding the title defective (\href{https://supreme.justia.com/cases/federal/us/14/304/}{14 U.S. at 353--60}).

\subsection{Concurring Opinion (Justice Johnson)}

Justice William Johnson concurred in the judgment but sharply distanced himself from Story's nationalist constitutional theory and mandatory-vesting dicta. Viewing the Constitution as a ``tripartite'' compact among the people, the States, and the United States (\href{https://supreme.justia.com/cases/federal/us/14/304/}{14 U.S. at 373}), Johnson rejected Story's claim that ``shall extend'' is imperative or that the omission of ``all'' creates two tiers of Article~III jurisdiction (\href{https://supreme.justia.com/cases/federal/us/14/304/}{14 U.S. at 374--75}). Instead, Johnson reasoned that Congress had discretion to assume or leave federal-question jurisdiction in state courts, and that Congress's greater power to exclude state courts altogether included the lesser power to permit concurrent state adjudication subject to federal revision (\href{https://supreme.justia.com/cases/federal/us/14/304/}{14 U.S. at 375--76}).

Crucially, Johnson emphasized comity and institutional restraint. Declaring ``I rejoice that Virginia has resisted'' insofar as resistance underscored the vitality of state judicial independence, Johnson argued that the Supreme Court exercises appellate power ``over persons and cases'' rather than compulsory \textit{in personam} control over state courts as subordinate bodies (\href{https://supreme.justia.com/cases/federal/us/14/304/}{14 U.S. at 362--63, 376, 379--82}). He construed Section~25's authorization for the Supreme Court to proceed directly to final decision and award execution after a prior remand as Congress's deliberate device to avoid issuing compulsory writs of mandamus to state judges (\href{https://supreme.justia.com/cases/federal/us/14/304/}{14 U.S. at 365--67}).

\subsection{Dissents}

There were no dissenting opinions.

\section{Examples: Future Applications}

\begin{enumerate}
    \item \textbf{Same-Side Application (State Criminal Conviction Raising a Federal Defense):} Suppose a state trial court convicts a defendant under a state lottery statute, rejecting the defendant's claim that a federal statute authorized the ticket sales, and the state's highest court affirms. Under \textit{Martin}, the Supreme Court has Article~III appellate jurisdiction to review the state criminal judgment because Article~III attaches to the federal question in the \textit{case}, not to whether the proceeding is civil or criminal or originated in state court---the exact extension the Court adopted five years later in \textit{\href{https://supreme.justia.com/cases/federal/us/19/264/}{Cohens v. Virginia}}, 19 U.S. (6 Wheat.) 264 (1821).

    \item \textbf{Same-Side Application (State Court Refusal to Enforce a Federal Supremacy Mandate):} Suppose a state supreme court invalidates a federal preemption defense in a tort suit, the U.S. Supreme Court reverses and remands, and the state supreme court issues an order declaring the federal statute unconstitutional and refusing to enter judgment for the defendant. Just as in \textit{Martin}, the state court's defiant order is a reviewable final judgment denying federal statutory validity, and the Supreme Court may reverse directly and affirm the proper judgment without needing to issue a writ of mandamus to state judges.

    \item \textbf{Opposite-Side Application (Pure State-Law Judgment Resting on an Adequate and Independent State Ground):} Suppose a state supreme court rules for a landlord in a contract dispute on two grounds: first, that a federal housing regulation does not apply; and second, that the tenant independently waived all defenses under a well-settled, neutral state rule of civil procedure. Unlike \textit{Martin}---where the state-law title question was inextricably antecedent to whether the Jay Treaty protected Fairfax's estate---the state procedural ground here independently supports the judgment regardless of how the federal issue is resolved, so the Supreme Court lacks jurisdiction to review or overturn the state court's interpretation of state law (\href{https://www.fjc.gov/history/cases/cases-that-shaped-the-federal-courts/martin-v-hunter-s-lessee}{FJC}; see also \textit{\href{https://supreme.justia.com/cases/federal/us/355/313/}{Staub v. City of Baxley}}, 355 U.S. 313 (1958)).

    \item \textbf{Opposite-Side Application (State Court Reviewing or Enjoining a Federal Court Judgment):} Suppose a federal district court enters a final judgment enforcing a federal treaty right, and the losing party petitions a state supreme court for a writ of error or habeas order setting aside the federal judgment on the theory that state and federal judiciaries are co-equal sovereigns. \textit{Martin}'s holding runs strictly one way under the Supremacy Clause: state courts have no reciprocal appellate or supervisory authority to overturn federal court judgments (\href{https://www.fjc.gov/history/cases/cases-that-shaped-the-federal-courts/martin-v-hunter-s-lessee}{FJC}, noting \textit{Ableman v. Booth} (1859)).

    \item \textbf{Fence-Sitter (Congressional Stripping of Supreme Court Appellate Review Over a Class of State-Court Federal Constitutional Claims):} Suppose Congress invokes Article~III's Exceptions Clause to strip both lower federal courts and the U.S. Supreme Court of jurisdiction over all constitutional challenges to state abortion or religious-exercise statutes, leaving final adjudication exclusively to state supreme courts. Under \textit{Martin}'s narrow holding (upholding Congress's \textit{grant} of jurisdiction in Section~25) and Justice Johnson's concurrence, Congress arguably possesses broad Exceptions Clause authority to leave such cases in state courts; yet under Justice Story's mandatory-vesting and two-tier dicta---later championed by \href{https://texaslawreview.org/congressional-power-to-strip-state-courts-of-jurisdiction/}{Akhil Reed Amar}---vesting zero federal judicial power (original or appellate) in ``all cases'' arising under the Constitution would violate Article~III's imperative text and destroy the uniformity \textit{Martin} deemed essential.
\end{enumerate}

\section{Critique}

\textbf{Originalist and Textualist Perspectives:} Contemporary states'-rights jurists led by Virginia Judge Spencer Roane and Judge William Cabell attacked Section~25 on strict textual and structural grounds, noting that Article~III never mentions ``state courts'' and arguing that subjecting sovereign state tribunals to federal writs of error impermissibly ``commandeered'' state judiciaries and converted them into inferior federal courts (\href{https://chicagounbound.uchicago.edu/cgi/viewcontent.cgi?article=1427&context=journal_articles}{Chicago Unbound}; \href{https://www.cambridge.org/core/books/monitoring-american-federalism/state-challenges-to-the-supreme-courts-control-over-constitutional-interpretation/85539C2F1811467C31EF62DBD350C7D8}{Cambridge}). Modern originalists and textualists, however, overwhelmingly vindicate \textit{Martin}'s core holding based on the Supremacy Clause, ratification-era assurances (including \textit{Federalist No.~82}), and the First Congress's enactment of Section~25. At the same time, textualists sharply divide over Story's \textit{dicta}. Neo-Federalist textualist \href{https://openyls.law.yale.edu/bitstreams/f4002d33-80ce-42a7-be16-e38c2f5930f0/download}{Akhil Reed Amar} embraces and refines Story's ``two-tier'' reading of ``shall be vested'' and ``all Cases'' to argue that Congress must vest final federal-court authority (original or appellate) over every federal-question, admiralty, and ambassadorial case. By contrast, critics such as Martin Redish and Daniel Meltzer (discussed in \href{https://openyls.law.yale.edu/bitstreams/f4002d33-80ce-42a7-be16-e38c2f5930f0/download}{Amar's reply}) argue that Story's mandatory-vesting dicta clashes with the Madisonian Compromise of 1787, which deliberately left the creation of lower federal courts to Congress's discretion, while \href{https://wwws.law.northwestern.edu/faculty/fulltime/pfander/stateinferiority.pdf}{James Pfander} offers an alternative textual synthesis grounded in Congress's Article~I power to ``constitute Tribunals inferior to the supreme Court.''

\textbf{Progressive, Doctrinal, and Institutional Perspectives:} Doctrinal and progressive scholars celebrate \textit{Martin} as the structural linchpin of national constitutional rights: without Supreme Court review of state judgments, the Bill of Rights, Reconstruction Amendments, and federal civil-rights statutes could be nullified state by state whenever local majorities resisted national law. Yet scholars also note a subtle tension in \textit{Martin}'s treatment of state law. Story's willingness to re-examine Virginia property law to prevent state-court evasion of the Jay Treaty was later cabined by the Court's ``adequate and independent state grounds'' doctrine (starting with \textit{Murdock v. Memphis} (1875) and \textit{Eustis v. Bolles} (1893); see \href{https://www.fjc.gov/history/cases/cases-that-shaped-the-federal-courts/martin-v-hunter-s-lessee}{FJC}), though the Court still polices state procedural grounds that are novel or manipulated to evade federal rights (\textit{\href{https://supreme.justia.com/cases/federal/us/355/313/}{Staub v. City of Baxley}}). Critics of judicial supremacy also observe that Story's opinion helped consolidate an exclusive judicial monopoly over constitutional interpretation, sidelining state legislatures as monitors of federalism (\href{https://www.cambridge.org/core/books/monitoring-american-federalism/state-challenges-to-the-supreme-courts-control-over-constitutional-interpretation/85539C2F1811467C31EF62DBD350C7D8}{Cambridge}).

\textbf{Competing Constitutional Values:} \textit{Martin} mediates a foundational clash between \textit{national uniformity and federal supremacy} (ensuring that one Constitution and one body of federal treaties and statutes govern equally across every state) and \textit{dual sovereignty and state judicial autonomy} (preserving state courts as independent organs of self-governing states rather than mere administrative adjuncts of Washington).

\section{Key Quotations}

\begin{quote}
``The Constitution of the United States was ordained and established not by the States in their sovereign capacities, but emphatically, as the preamble of the Constitution declares, by `the people of the United States.'''
\end{quote}
\noindent --- Justice Joseph Story, Opinion of the Court, \href{https://supreme.justia.com/cases/federal/us/14/304/}{14 U.S. (1 Wheat.) at 324} (also at \href{https://www.courtlistener.com/opinion/85160/martin-v-hunters-lessee/}{CourtListener}).

\begin{quote}
``It is the case, then, and not the court, that gives the jurisdiction. If the judicial power extends to the case, it will be in vain to search in the letter of the Constitution for any qualification as to the tribunal where it depends.''
\end{quote}
\noindent --- Justice Joseph Story, Opinion of the Court, \href{https://supreme.justia.com/cases/federal/us/14/304/}{14 U.S. (1 Wheat.) at 338} (also at \href{https://www.courtlistener.com/opinion/85160/martin-v-hunters-lessee/}{CourtListener}).

\begin{quote}
``That motive is the importance, and even necessity, of uniformity of decisions throughout the whole United States upon all subjects within the purview of the Constitution. Judges of equal learning and integrity in different States might differently interpret a statute or a treaty of the United States, or even the Constitution itself; if there were no revising authority to control these jarring and discordant judgments and harmonize them into uniformity, the laws, the treaties, and the Constitution of the United States would be different in different States, and might perhaps never have precisely the same construction, obligation, or efficacy in any two States.''
\end{quote}
\noindent --- Justice Joseph Story, Opinion of the Court, \href{https://supreme.justia.com/cases/federal/us/14/304/}{14 U.S. (1 Wheat.) at 347--48} (also at \href{https://www.courtlistener.com/opinion/85160/martin-v-hunters-lessee/}{CourtListener}).

\begin{quote}
``On the other hand, so firmly am I persuaded that the American people can no longer enjoy the blessings of a free government whenever the State sovereignties shall be prostrated at the feet of the General Government, nor the proud consciousness of equality and security any longer than the independence of judicial power shall be maintained consecrated and intangible, that I could borrow the language of a celebrated orator and exclaim, `I rejoice that Virginia has resisted.'''
\end{quote}
\noindent --- Justice William Johnson, Concurring Opinion, \href{https://supreme.justia.com/cases/federal/us/14/304/}{14 U.S. (1 Wheat.) at 363} (also at \href{https://www.courtlistener.com/opinion/85160/martin-v-hunters-lessee/}{CourtListener}).

\end{document}
